\documentclass[12pt]{article}
\usepackage[utf8]{inputenc}
\usepackage[T1]{fontenc}
\usepackage{hyperref}
\title{Rawo Language Resources for K-12 Curriculum Development}
\author{Compiled by Claude Sonnet 5.5 (claude-sonnet-5-5) for Andr\'as Kornai}
\date{}
\begin{document}
\maketitle

\section{Introduction}
Rawo (ISO 639-3 \texttt{rwa}; also called Poko or Poko-Rawo) is a Papuan language of the Skou family. It is spoken on the north coast of Sandaun Province, Papua New Guinea, near the village of Leitre. The sources give about 640 speakers (2003). Joshua Project places the people in East Sepik Province, which conflicts with the linguistic sources. Very little educational material turned up in search results.

\section{Primary Educational Resources}
No primers, textbooks or school programmes were found.

\section{Linguistic and Reference Materials}
\begin{itemize}
\item Wikipedia, ``Rawo language'': overview with classification and speaker count. \url{https://en.wikipedia.org/wiki/Rawo_language}
\item Verbix, Rawo language page. \url{https://docs.verbix.com/Languages/Rawo}
\item WALS language entry (code \texttt{pok}). \url{https://wals.info/languoid/lect/wals_code_pok}
\item Wikipedia, ``Skou languages'' (mirror copy): background on the family. \url{https://static.hlt.bme.hu/semantics/external/pages/jokuc/en.wikipedia.org/wiki/Skou_languages.html}
\end{itemize}
These are reference pages only. No dictionary or full grammar was located.

\section{Religious and Community Resources}
\begin{itemize}
\item Joshua Project, ``Rawo in Papua New Guinea'': people-group profile. \url{https://joshuaproject.net/people_groups/14530}
\item Joshua Project language page for \texttt{rwa}. \url{https://joshuaproject.net/languages/rwa}
\end{itemize}

\section{Digital and Audio Resources}
No audio, corpus or software resources were found.

\section{Curriculum Development Context}
The language is not known to be taught in schools. With roughly 640 speakers, and Tok Pisin and English dominant in PNG education, any curriculum would need community-led documentation first. Primary-level vernacular preschool or bridging materials would be the realistic starting point.
\end{document}